\documentclass{article}
\usepackage{amsmath,amssymb}
\usepackage{xurl}
\usepackage[hidelinks]{hyperref}
% Shared commands for notes in docs/paper-gaps/.

\DeclareUnicodeCharacter{2080}{\ifmmode{}_{0}\else\textsubscript{0}\fi}
\DeclareUnicodeCharacter{2081}{\ifmmode{}_{1}\else\textsubscript{1}\fi}
\DeclareUnicodeCharacter{2082}{\ifmmode{}_{2}\else\textsubscript{2}\fi}
\DeclareUnicodeCharacter{2099}{\ifmmode{}_{n}\else\textsubscript{n}\fi}

\newcommand{\Lean}{\textsc{Lean}}
\ExplSyntaxOn
\NewDocumentCommand{\leanid}{m}{
  \ifmmode
    \textsf{\detokenize{#1}}
  \else
    \group_begin:
    \urlstyle{sf}
    \tl_set:Nn \l_tmpa_tl {#1}
    \tl_replace_all:Nnn \l_tmpa_tl {\_} {_}
    \exp_args:NV \path \l_tmpa_tl
    \group_end:
  \fi
}
\ExplSyntaxOff
\providecommand{\tr}{\operatorname{tr}}
\newcommand{\tnleanIssueBase}{https://github.com/LionSR/TNLean/issues/}
\newcommand{\tnleanPrBase}{https://github.com/LionSR/TNLean/pull/}
\newcommand{\tnleanDocBase}{https://sirui-lu.com/TNLean/docs/}
\newcommand{\ghissue}[1]{\href{\tnleanIssueBase#1}{GitHub issue \##1}}
\newcommand{\ghpr}[1]{\href{\tnleanPrBase#1}{PR \##1}}
\newcommand{\doclink}[1]{\href{\tnleanDocBase#1.html}{\path{#1}}}

% Dirac notation and the identity operator, as in blueprint/src/macros/common.tex.
% \providecommand keeps note-local definitions authoritative.
\usepackage{dsfont}
\providecommand{\ket}[1]{|#1\rangle}
\providecommand{\bra}[1]{\langle #1|}
\providecommand{\braket}[2]{\langle #1 | #2 \rangle}
\providecommand{\ketbra}[2]{|#1\rangle\!\langle #2|}
\providecommand{\Id}{\mathds{1}}
\providecommand{\one}{\mathds{1}}
\providecommand{\C}{\mathbb{C}}
\providecommand{\MN}[1]{M_{#1}(\C)}
\providecommand{\MD}{M_D(\C)}

% External referencing for paper sources.  The build helper writes sanitized
% auxiliary files under build/paper-aux/ and excludes duplicated source labels.
\usepackage{xr}

% Import a generated paper auxiliary file with a caller-supplied label prefix.
% The candidate paths support builds from docs/paper-gaps/, the repository root,
% and build/paper-aux/ itself.
\newcommand{\importpaperaux}[2]{%
  \IfFileExists{../../build/paper-aux/#2.aux}{%
    \externaldocument[#1]{../../build/paper-aux/#2}%
  }{%
    \IfFileExists{build/paper-aux/#2.aux}{%
      \externaldocument[#1]{build/paper-aux/#2}%
    }{%
      \IfFileExists{#2.aux}{%
        \externaldocument[#1]{#2}%
      }{}%
    }%
  }%
}

% General external reference macros
\newcommand{\paperref}[2]{\ref{#1-#2}}
\newcommand{\papereqref}[2]{(\ref{#1-#2})}

% CPSV16 source (1606.00608 / MPDO-22-12-17-2.tex) external document and shortcuts
\importpaperaux{cpsv16-}{1606.00608/MPDO-22-12-17-2-tnlean}
\newcommand{\cpsvref}[1]{\paperref{cpsv16}{#1}}
\newcommand{\cpsveqref}[1]{\papereqref{cpsv16}{#1}}

% Verdict marker: \gapnote{<kind>}{<status>}. Kind is the policy
% classification (clarification, local-correction, scope-restriction,
% unfaithful, false-source, open-gap); status is open, wip (elimination
% underway), resolved, or historical. Machine-read by the site index;
% prints a verdict line.
\newcommand{\gapnote}[2]{\par\noindent\textbf{Verdict:} #1 (#2).\par}

\title{Cocycle Representatives in the Stabilizer Extension Calculation}
\author{}
\date{2026-09-04}
\begin{document}
\maketitle
\gapnote{scope-restriction}{open}

Proposition \texttt{prop:BI\_psi} of
Ref.~\cite{FrancoRubio2025MPUGauging}, lines 7042--7064, writes
$\Psi\in H^2(G,U(1))$ and then evaluates $\Psi(g,h)$.
Evaluation is defined for a cocycle representative, not for a cohomology class.
At the class level the hypothesis is
$[\psi]\in\operatorname{im}(\operatorname{res}^G_H)$.
At the representative level the displayed calculation uses a genuine
cocycle $\Psi$ satisfying $\Psi|_{H\times H}=\psi$ after compatible
coboundary choices.

The present scalar result assumes exactly this latter equality and proves
\begin{align}
 \gamma_g^x&=\frac{\Psi(k_{gx}^{-1},gk_x)\Psi(g,k_x)}
 {\Psi(e,e)\Psi(k_x,k_x^{-1})},\\
 \psi(t_k(g_1,g_2x),t_k(g_2,x))
 &=\Psi(g_1,g_2)\frac{\gamma_{g_1g_2}^x}
 {\gamma_{g_1}^{g_2x}\gamma_{g_2}^x}.
\end{align}
Here $t_k(g,x)=k_{gx}^{-1}gk_x$, with the corrected index in the first
transition factor documented in the companion representative-index note.
The printed formula for $\gamma$ itself requires no correction, including
its denominator involving $k_x$ rather than $k_{gx}$.
The inverse gauges trivialize the induced scalar symbols.
No normalization of $\Psi$, finiteness, or normality of $H$ is assumed.
The result is proved for $\mathbb C^\times$ coefficients, and hence applies
to unit-circle-valued cocycles viewed as nonzero complex scalars.

\textbf{Remaining scope.}
This is only the forward scalar representative calculation. It does not
construct fusion or action tensors, prove the converse, or identify the
class-level restriction image. Issue \#7361 remains open for the full
proposition. Its completion must justify compatible coboundary choices,
prove the converse and class-level equivalence, and connect the scalar
calculation to the tensor hypotheses. No quotient group structure on
$G/H$ is used or required.

\bibliographystyle{alpha}
\bibliography{references}
\end{document}
